%% ============================================================
%% 分片 10：原文第 4 节「Compactification And The Geometry Of Hitchin's Moduli Space」（原文 glshort.tex 起于第 2325 行）
%% 由 tools/harvmac2tex.py 从 arXiv 源码（hep-th/0604151v3）自动生成。
%% 只译英文散文；公式、编号、\label、\cite 一律不改。
%% ============================================================

\subsection{$\MH$ 的复结构}

\label{sec:complexstr}

我们下一步的任务是更具体地理解 $\MH$ 上的复结构。

\par\medskip\emph{用全纯数据描述}

让我们从复结构 $I$ 的角度来考察 $\EUW$ 的超凯勒商。组合
$\nu_I=\mu_J+i\mu_K$ 在这个复结构下是全纯的：
% 式 (4.20)
\setcounter{equation}{19}
\begin{equation}\label{eq:uttup}
\nu_I=-{1\over \pi}\int_C|d^2z|\,\Tr\,\epsilon\,D_{\bar
z}\phi_z.
\end{equation}
 $\nu_I$ 在复结构 $I$ 下全纯，因为它是从 $\Omega_I$ 导出的 moment map，而
后者在复结构 $I$ 下是 $(2,0)$ 型的（具体地说，$\nu_I$ 之所以全纯，是因为它只依赖于 $A_{\bar z}$ 与 $\phi_z$）。
在这类问题中，把全纯 moment map $\nu_I$ 的零点与实 moment map $\mu_I$ 的零点分开考虑，往往是很方便的 \cite{ref61}。

黎曼面 $C$ 上光滑 $G$-丛 $E$ 上的任何联络 $A$ 都自动把 $E$ 变成全纯 $G$-丛
（我们仍用同一个名字记它，但愿这不会引起混淆）。只需把 $\bar\partial$ 算子定义为
$\bar D=d\bar z\,D_{\bar z}$；在复维数为一的情形，它平凡地满足 $\bar D^2=0$。全纯
moment map $\nu_I$ 为零告诉我们 $\bar D\phi=0$；换言之，
$\varphi=\phi_z\,dz$ 是 $K_C\otimes {\rm
ad}(E)$ 的一个全纯截面，其中 $K_C$ 是 $C$ 的典范线丛。
换个说法，$\varphi$ 是一个取值于伴随表示、也就是取值于 $E$ 的伴随丛的全纯
1-形式。我们把由 $G$-丛 $E$ 与 $K_C\otimes {\rm ad}(E)$ 的全纯截面 $\varphi$ 组成的
对 $(E,\varphi)$ 称为 Hitchin 对。$\varphi$ 常被称为 Higgs 场，而赋以 $\varphi$
之选择的丛 $E$ 也称为 Higgs 丛。

要得到模空间 $\MH$，我们还必须把实 moment map $\mu_I$ 置零，并模去取值于 $G$ 的
规范变换所构成的群。然而，正如 \cite{ref20} 所利用的（这也是模问题中常有的情形），
理解这些合在一起的步骤有一个更简单的办法。全纯 moment map $\nu_I$ 的零点空间 $\EUW_0$
不仅容许通常取值于 $G$ 的规范变换作用，还容许取值于 $G$ 的复化 $G_{\Bbb{C}}$ 的
规范变换作用。这一点从 $\nu_I$ 的全纯性看得很清楚。于是我们可以对 $\EUW_0$
施行以下两种操作之一：

\emph{(I)} 限制到 $\EUW_0$ 中满足 $\mu_I=0$ 的子空间，再模去取值于 $G$ 的规范变换。

\emph{(II)} 用取值于 $G_{\Bbb{C}}$ 的规范变换去除 $\EUW_0$。

操作 \emph{(I)} 给出所要的模空间
$\MH$，但操作 \emph{(II)} 容易理解得多，且与它几乎重合。
（两者几乎重合的原因是：几乎每条复规范变换轨道都包含唯一一条满足 $\mu_I=0$ 的通常规范变换轨道。）
理解操作 \emph{(I)} 最容易的办法，常常是先理解操作 \emph{(II)}，再理解它与操作 \emph{(I)} 的关系。

在当前的例子中，操作 \emph{(II)} 的结果很容易描述。它的意思是：我们不关心 $A_{\bar z}, $ $\phi_z$ 的具体取法，
而只关心由 $A_{\bar z}$ 决定的全纯 $G$-丛 $E$，以及与之相关联的
$K_C\otimes {\ad}(E)$ 的全纯截面 $\varphi$。
因此，在当前的例子中，操作 \emph{(II)} 给出的是 Hitchin 对 $(E,\varphi)$ 的等价类的集合。

这个集合在许多用途上都是 $\MH$ 的极好近似；
例如当 $C$ 的亏格至少为 2 时，两者只相差余维相当高的集合。更准确地说（本文大部分
内容并不依赖于这些细节），我们需要稳定性这一概念。
对 $G=SU(2)$，我们把 $E$ 解释为 $C$ 上秩为二的复全纯向量丛，把 $\varphi$ 解释为全纯映射 $\varphi:E\to E\otimes K_C$。
若一个线丛 ${\cal L}\subset E$ 满足 $\varphi({\cal L})\subset {\cal L}\otimes K_C$，则称它
为 $\varphi$-不变的。若每个 $\varphi$-不变的线丛 ${\cal L}\subset E$ 的第一 Chern 类都是负的，则称
Hitchin 对 $(E,\varphi)$ 为稳定的。若每个这样的 ${\cal L}$ 的第一 Chern 类都非正，则称它
为半稳定的。对一般的 $G$，必须考虑把 $E$ 的结构群约化到 $G_{\Bbb{C}}$ 的极大抛物子群 $P$ 的 $\varphi$-不变约化。带这种约化的丛 $E$ 有一个自然的第一 Chern 类
（因为 $P$ 带有一个特殊的 $U(1)$）；若对每个这样的约化，第一 Chern 类都是负的（或非正的），则称对
$(E,\varphi)$ 为稳定（或半稳定）的。半稳定但不稳定的对称为严格半稳定的。

稳定性是一个温和的条件，意思是：例如若
$(E,\varphi)$ 是稳定的，则每个邻近的对也是稳定的。此外，不稳定的对具有很高的余维（如果 $C$ 的亏格大于 1）。

现在我们可以回到比较操作 \emph{(I)} 与 \emph{(II)} 的问题，或者等价地说，回到用全纯语言描述 Hitchin 方程解的模空间 $\MH$ 的问题。
\cite{ref20} 证明的结果是：$\MH$ 等同于「稳定对 $(E,\varphi)$ 的模空间」，即稳定 Hitchin 对。要得到这个模空间，我们扔掉不稳定的 Hitchin 对，对半稳定的那些施加某种等价关系，然后模去复值规范变换。这个对操作 \emph{(II)} 稍作修改的版本——称之为操作 \emph{(II)$'$}——与操作 \emph{(I)} 恰好一致。

总之，$\MH$ 参数化稳定的 Hitchin 对 $(E,\varphi)$，以及严格半稳定 Hitchin 对的某些等价类。稳定对对应于 $\MH$ 中的光滑点，也对应于（对某些 $G$）某些轨形奇点（orbifold 奇点）。严格半稳定的对一般会导致 $\MH$ 中比轨形奇点更糟糕的奇点。

\par\medskip\emph{复结构 $J$ 的类比}

以上这些在复结构 $J$ 中都有类比。复结构 $J$ 中的全纯 moment map 是 $\nu_J=\mu_K+i\mu_I$：
% 式 (4.21)
\setcounter{equation}{20}
\begin{equation}\label{eq:nuJ}
\nu_J =-{i\over 2\pi}\int_C\,\Tr\,\epsilon{\cal F}
\end{equation}
 为
构造 $\MH$，我们首先要求 $\nu_J$ 为零。$\nu_J$ 的一个零点不过是一对 $A,\phi$，使得曲率
${\cal F}$（属于取值于 $G_{\Bbb{C}}$ 的联络 $\CA=A+i\phi$）等于零。因此，如果我们用取值于 $G_{\Bbb{C}}$ 的
规范变换群去除 $\nu_J$ 的零点集，就恰好得到全部同态 $\vartheta:\pi_1(C)\to G_{\Bbb{C}}$ 在共轭下的集合 ${\cal Y}_0$。这就是操作 \emph{(II)}。

相反，要得到 $\MH$，我们应当施行操作 \emph{(I)}。换句话说，我们应当置 $\mu_J=0$，即施加条件 $D^*\phi=0$，并且只模去取值于 $G$ 的规范变换。
与复结构 $I$ 的情形一样，在比较这两种操作时，稳定性这一概念会介入。
对 $G=SU(2)$，若同态 $\vartheta:\pi_1(C)\to G_{\Bbb{C}}$ 的单值不能同时化归为三角形式 % 式 (4.22)
\setcounter{equation}{21}
\begin{equation}\label{eq:yto}
\left(\begin{matrix} \alpha &
\beta \\ 0 & \alpha^{-1}\end{matrix}\right).
\end{equation}
   则称该同态为稳定的。若 $\vartheta$ 有这样的化归，则称它为严格半稳定的。（在复结构
$J$ 中，没有不稳定表示的概念。）若两个严格半稳定表示具有相同的对角单值元素，即相同的 $\alpha$，则认为它们等价。
注意每个这样的等价类都含一个 $\beta=0$ 的特殊代表元。

对任何 $G$，把单值放进三角形式的类比，是把它们共轭进 $G$ 的一个抛物子群 $P$ 中。若一个表示不能被这样共轭，则称它为稳定的，
否则称为严格半稳定的。严格半稳定表示上的等价关系对任何 $G$ 都有自然的类比
（两个都约化到 $P$ 的严格半稳定表示等价，如果这两个平坦 $P$-丛投影到 $P$ 的极大约化子群后变得等价）。

Corlette \cite{ref58} 与 Donaldson 的一个定理（见 \cite{ref20} 的附录）说：$\MH$ 在复结构 $J$ 中就是
参数化稳定同态 $\vartheta:\pi_1(C)\to G_{\Bbb{C}}$ 以及半稳定同态之等价类的模空间 ${{\cal Y}}$。

\par\medskip\emph{$I_w$ 的类比}

我们刚才所说的，在每个复结构 $I_w$（$w\not= 0,\infty$）中都有直接的类比。复结构 $I_w$ 中的全纯
变量是 $A_{\bar z}-w\phi_{\bar z}$、$A_z+w^{-1}\phi_z$。
Hitchin 方程中的两个方程合起来给出全纯条件 $[D_{\bar z}-w\phi_{\bar
z},D_z+w^{-1}\phi_z]=0$，第三个是一个 moment map 条件。
该全纯条件说明：分量为 $(A_{\bar z}-w\phi_{\bar z},A_z+w^{-1}\phi_z)$ 的复值联络是平坦的。
把 moment map 置零并模去规范变换，得到的模空间与我们在复结构 $J$ 中找到的 ${{\cal Y}}$ 相同。
因此所有复结构 $I_w$（$w\not=0,\infty$）都是等价的。

事实上，在复结构 $I=I_{w=0}$ 中，我们可以定义一个对称群 ${\cal U}\cong \Bbb{C}^*$，它作用在
Higgs 丛上为 $(E,\varphi)\to (E,\lambda\varphi)$，
$\lambda\in\Bbb{C}^*$。（稳定性条件在这个变换下不变。）当 $|\lambda|=1$ 时，它约化为第~\ref{sec:mhhyper}~节中描述的群
${\cal U}_1$，后者是 $\MH$ 的超凯勒度量的一个对称群。当 $|\lambda|\not=1$ 时，
我们得不到 $\MH$ 的超凯勒度量的对称，但正如 \cite{ref20} 第 107-8 页所示，
我们得到 $\MH$ 的一个微分同胚群，它保持复结构 $I$，并通过 % 式 (4.23)
\setcounter{equation}{22}
\begin{equation}\label{eq:trymodo}
I_w\to
I_{\lambda^{-1}w},
\end{equation}
 变换其他复结构，从而把式 \eqref{eq:zungo}\ 推广到
$|\lambda|\not=1$。

\subsection{Hitchin 纤维化}

\label{sec:hitchfib}

现在我们来更详细地讨论 $\MH$ 的几何。

若 $E$ 是稳定的 $G$-丛，则对 $(E,0)$ 是稳定的
Hitchin 对。于是我们得到 $C$ 上稳定 $G$-丛的模空间 $\M$ 到 Hitchin 模空间
$\MH$ 的自然嵌入。在复结构 $I$ 中，$\M$ 是 $\MH$ 的全纯子流形，因为它由方程
$\phi_z=0$ 定义，而这些方程在复结构 $I$ 中全纯。在复结构
$J$ 与 $K$ 中，$\M$ 不是全纯的。相反，它是 Lagrangian 的，因为 $\phi=0$ 蕴含 $\delta\phi=0$，从而
$\mu_J=\mu_K=0$。

若 $E$ 是稳定的，则对每个 $\varphi\in H^0(C,K_C\otimes\ad(E))$，Hitchin 对 $(E,\varphi)$ 都是稳定的。此外，${\EUM}$ 在点 $E$ 处的
切丛是 $H^1(C,\ad(E))$（它参数化 $E$ 的一阶形变），故由 Serre
对偶，$\M$ 在点 $E$ 处的余切空间是
$H^0(C,K_C\otimes\ad(E))$。于是 $\varphi$ 取遍这个
余切空间中的值，而所有满足 $E$ 稳定的 $(E,\varphi)$ 构成的空间就是余切丛
$T^*\M$。因此我们实际上得到
$T^*\M$ 到 $\MH$ 的嵌入。$\MH$ 在复结构 $I$ 中的全纯辛形式
$\Omega_I$ 限制到 $T^*\M$ 上，正好是它作为全纯余切丛的自然辛结构
（也就是说，对 $\M$ 上任意的局部全纯坐标 $q^\alpha$
，存在在 $T^*\M\to \M$ 的纤维上线性的全纯函数 $p_\alpha$，使得
$\Omega_I=\sum_\alpha dp_\alpha\wedge dq^\alpha$）。\footnote{这个断言成立，是因为类似
的断言在取超凯勒商之前的全部对 $(A,\phi)$ 构成的空间 $\EUW$ 上实际上成立。$\EUW$
可以理解为联络空间 ${\EUA}$ 的余切空间，$A$ 参数化 ${\EUA}$，而 $\phi$
参数化 ${\EUA}$ 的余切空间；$\Omega_I$ 就是 ${\EUW}$ 看作
$T^*{\EUA}$ 时的自然全纯辛形式。一般来说，这种类型的结构在取复辛商
（在现下的情形中，即置 $\nu_I=0$ 并模去取值于 $G_{\Bbb{C}}$ 的规范变换）时保留下来。}

$T^*\M$ 在 $\MH$ 中的像并不是整个 $\MH$，因为 Hitchin 对 $(E,\varphi)$ 可能在 $E$ 不稳定时仍然稳定。
然而，那些 $E$ 不稳定的稳定 Hitchin 对 $(E,\varphi)$ 构成一个余维非常高的集合。扔掉
这个集合后，$\MH$ 同构于 $T^*\M$，并且有一个通过忘掉 $\varphi$ 而得到的到
$\M$ 的自然映射。尽管这个映射只在去掉一个余维非常高的集合之后才有定义，它却
极其有用。我们称它为 Hitchin 的第二纤维化（第一个纤维化是我们马上就要引入的另一个映射）。

\par\medskip\emph{Hitchin 纤维化}

在复结构 $I$ 中，除了把 $\varphi$ 映为零之外，另一个自然的操作
是计算 $\varphi$ 的规范不变多项式。对 $G=SU(2)$，这简单意味着我们考虑
二次 Casimir 算子 $w=\Tr\,\varphi^2$。由于
$\bar D\varphi=0$，我们有 $\bar\partial w=0$，所以 $w$ 是一个全纯二次微分，取值于
${\EUBB}=H^0(C,K_C^2)\cong \Bbb{C}^{3g-3}$。最常被称为 Hitchin 纤维化的映射，是把对 $(E,\varphi)$ 映到
$w=\Tr\,\varphi^2$ 所得的映射 $\pi:\MH\to {\EUBB}$。

对任何 $G$，Hitchin 纤维化的定义类似，只是人们用 $\varphi$ 的全部独立
Casimir（即全部独立的 $G$-不变多项式）来刻画 $\varphi$，而不只是二次的那个。例如
对 $G=SU(N)$，我们定义 $w_n=\Tr\,\varphi^n$
（$n=2,\dots,N$），并令 ${\EUBB}=\oplus_{n=2}^NH^0(C,K_C^n)$。
于是 Hitchin 纤维化就是把 $(E,\varphi)$ 映到点 $(w_2,w_3,\dots,w_n)\in {\EUBB}$ 的映射。对任何 $G$，人们考虑
的则是代替 $\Tr\,\varphi^n$ 的适当独立齐次 $G$-不变多项式 ${\cal P}_i$。这些多项式的个数
等于 $G$ 的秩 $r$，它们的次数 $d_i$ 满足恒等式 % 式 (4.24)
\setcounter{equation}{23}
\begin{equation}\label{eq:yty}
\sum_{i=1}^r(2d_i-1)={\rm
dim}(G).
\end{equation}
 例如对 $G=SU(N)$，$d_i$ 是 $2,3,\dots,N$，
于是 $\sum_i(2d_i-1)=N^2-1={\rm dim}(G)$。最后，一般地我们定义 ${\EUBB}=\oplus_i H^0(C,K_C^{d_i})$。

由于 $\dim \,H^0(C,K_C^n)=(2n-1)(g-1)$，由 \eqref{eq:yty}\
可知 ${\EUBB}$ 的复维数是 $(g-1){\rm dim}(G)$，
它等于 $\M$ 的维数，也等于
$\M_H$ 维数的一半。Hitchin 纤维化 $\pi:\MH\to {\EUBB}$ 是满射，这一点我们马上就会讨论。由于 Hitchin 纤维化的底空间
${\EUBB}$（即映射 $\pi:\MH\to {\EUBB}$ 的靶空间）的维数是 $\MH$ 维数的一半，可知 $\pi$ 的一般纤维 $\CMF$ 的维数也是 $\MH$ 维数的一半，
从而等于 ${\EUBB}$ 的维数：
% 式 (4.25)
\setcounter{equation}{24}
\begin{equation}\label{eq:yrof}
\dim\,\CMF=\dim\,{\EUBB}={1\over
2}\dim\,\MH=(g-1)\dim\,G.
\end{equation}

让我们定性地解释为什么 Hitchin 纤维化是满射。例如取 $G=SU(2)$。取一个稳定的 $SU(2)$
丛 $E$。考虑方程 $\Tr\,\varphi^2=w$，其中
$\varphi $ 在 $(3g-3)$ 维空间
$H^0(C,K_C\otimes \ad(E))$ 中变化，而 $w$ 是
$(3g-3)$ 维空间 ${\EUBB}=H^0(C,K_C^2)$ 中的固定元素。这是
$3g-3$ 个复变量的 $3g-3$ 个二次方程组成的方程组。
解的一般个数是 $2^{3g-3}$。对其他 $G$ 也可以作类似的计数。

\par\medskip\emph{完全可积性}

\cite{ref62}

 我们现在想解释 Hitchin 的一个主要结果 ：$\MH$ 在复结构 $I$ 中是一个完全可积的哈密顿系统。我们可以找到 ${1\over 2}{\rm
dim}\MH$ 个 $\MH$ 上的函数 $H_a$，它们在复结构
$I$ 中全纯、代数无关，并且关于由全纯辛形式 $\Omega_I$ 得到的 Poisson 括号两两
Poisson 交换。\footnote{读者可能不习惯这种全纯意义上的完全可积系统。
从这样的系统可以提取通常实意义上的完全可积哈密顿系统（而且有趣且重要的构造正以这种方式出现；例子见
\cite{ref63,ref64}）。我们选取 $C$ 使它能容许一个实结构——即一个反转其复结构的对合。这在 $\M$ 与 $\MH$ 上诱导出实结构。
通过特殊化到 $\MH$ 中的一个实切片，就得到完全可积的哈密顿系统，其中相空间坐标、辛结构与交换的哈密顿量全都是实的。}

事实上，我们可以取 $H_a$ 为 ${\EUBB}$ 上的线性函数，因为 ${\EUBB}$ 的维数恰好等于所需的函数个数。
我们先对 $G=SU(2)$ 解释这一构造。首先取 $(3g-3)$ 维空间 $ H^1(C,T_C)$ 的一组基 $\alpha_a$
（$a=1,\dots,3g-3$），该空间与 $H^0(C, K_C^2)\cong {\EUBB}$ 对偶。（这里 $T_C$ 是
$C$ 的全纯切丛。）我们把 $\alpha_a$ 用取值于 $T_C$ 的 $(0,1)$-形式表示，并仍用同一个名字称呼它们，
然后定义 % 式 (4.26)
\setcounter{equation}{25}
\begin{equation}\label{eq:miro}
H_a=\int_C\alpha_a \wedge \Tr\,\varphi^2.
\end{equation}
 我们断言，这些函数关于全纯辛形式 $\Omega_I$ 两两 Poisson 交换。

一个自然的证明利用了下述事实：$H_a$ 的定义在取超凯勒商之前、在无穷维空间 $\EUW$ 上就有意义。
用 $\EUW$ 上的辛结构 $\Omega_I$ 来定义 Poisson 括号，$H_a$ 显然两两 Poisson 交换。因为在这些 Poisson 括号下（给定 $\Omega_I$ 的形式 \eqref{eq:ninto}\ ），$\varphi_z$ 与 $A_{\bar z}$ 共轭，
而与自身交换。但 $H_a$ 只是 $\varphi_z$ 的函数，不是 $A_{\bar z}$ 的函数。

这些函数 $H_a$ 可以限制到 $\nu_I=0$ 的轨迹上，
又因为它们在取值于 $G_{\Bbb{C}}$ 的规范变换下不变，它们下降到 $\MH$ 上的全纯函数。辛约化
的一个一般性质是：它把 Poisson 交换的函数映为 Poisson 交换的函数
（辛约化即把某个 moment map——这里是 $\nu_I$——置零，再模去相应的群——这里是取值于 $G_{\Bbb{C}}$ 的规范变换群）。
所以 $H_a$ 作为 $\MH$ 上的函数两两 Poisson 交换。它们的个数足够多，足以建立 $\MH$ 的完全可积性。

这个构造中有一点需要注意：作为 $\EUW$ 上的函数，$H_a$
依赖于代表上同调类的、取值于 $T$ 的 $(0,1)$-形式 $\alpha_a$ 的具体选取。
任何选取都可以，但我们必须作出选择。然而在限制并下降到 $\MH$ 之后，我们在 $\MH$ 上得到的函数只依赖于 $\alpha_a$ 的上同调类。事实上，一旦有 $\bar D\varphi=0$，从而
$\bar\partial\Tr\,\varphi^2=0$，$H_a$ 显然在
$\alpha_a\to\alpha_a+\bar\partial\epsilon_a$ 下不变。

Poisson 交换的函数 $H_a$ 在
$\MH$ 上生成彼此交换的流，这些流在复结构 $I$ 中全纯。
复环面容许交换的流，人们可能会猜测由 $H_a$ 生成的轨道至少一般地是复环面。这一点
由关于全纯辛结构与 Hitchin 纤维化的纤维紧性的普遍结果推出，也可以更直接地证明，即使用
谱覆盖技术 \cite{ref20,ref62}。这项技术在本课题的进一步发展中有
许多应用，将在别处解释 \cite{ref26}。在本文中，我们将在不使用谱覆盖构造的前提下尽量推进。

对任何 $G$，上述构造的类比是把 $\Tr\,\varphi^2$ 换成一个次数为 $d_i$ 的一般规范不变多项式 ${\cal
P}_i$。相应的交换哈密顿量取形式 $H_{\alpha,i}=\int_C \alpha \,{\cal P}_i(\varphi)$，其中
$\alpha\in H^1(C,K_C^{1-d_i})$。利用维数公式 \eqref{eq:yty} 作一个简单的维数计数即可说明，这些哈密顿量
在数目上恰好足以建立 $\MH$ 的完全可积性。利用 Hitchin 纤维化的纤维是紧的（从而全纯函数必是该纤维化底空间上的拉回）这一事实，可以证明这些交换哈密顿量生成
$\MH$ 上的全纯函数环。

完全可积性的一个简单而重要的推论是：Hitchin 纤维化的纤维是
$\Omega_I$ 这一全纯辛结构中的 Lagrangian 子流形，等价地也是实辛结构 $\omega_J$ 与
$\omega_K$ 中的 Lagrangian 子流形。事实上，该纤维化的一条纤维由方程 $H_k-h_k=0$ 定义，其中 $H_k$ 是交换的哈密顿量
而 $h_k$ 是复常数。一般来说，一组中间维数的 Poisson 交换函数（如
现下的 $H_k-h_k$）的零点集是 Lagrangian 的。

\begin{yizhu}[易错点]
稳定性的方向容易记反：对 $G=SU(2)$，稳定要求每个 $\varphi$-不变线丛 ${\cal L}\subset E$ 的第一 Chern 类为负（半稳定则要求非正）。这里的「负」是就 ${\cal L}$ 的第一 Chern 数即次数而言，不是就某个斜率的大小而言。对一般 $G$，要看的对象变成 $E$ 到极大抛物子群 $P$ 的 $\varphi$-不变约化，其第一 Chern 类由 $P$ 中那个特殊的 $U(1)$ 给出。
\end{yizhu}

\begin{yizhu}[含糊点拨]
为什么操作 (II) 与操作 (I) 几乎重合？原文只说「几乎每条复规范变换轨道都包含唯一一条满足 $\mu_I=0$ 的通常规范变换轨道」。按：这是 Kempf-Ness 型的论断——复群 $G_{\Bbb{C}}$ 的轨道与 $\mu_I^{-1}(0)$ 的交在稳定点处是单个 $G$-轨道，只有严格半稳定的轨道才会出现多解，所以本小节末尾要额外加上「半稳定对之间的等价关系」这一层。
\end{yizhu}

\begin{yizhu}[物理图像]
$T^*\M\hookrightarrow\MH$ 给出了 Hitchin 模空间的直观图像：在复结构 $I$ 下，$\MH$ 在余维很高的坏集之外就是稳定 $G$-丛模空间 $\M$ 的余切丛，而 Higgs 场 $\varphi$ 正是余切方向上的坐标。按：这正是几何 Langlands 把「$G$-丛上的 Hecke 修正」与「$\MH$ 中膜的形变」联系起来的几何基础（参见 Witten 2008 综述第 2 节）。
\end{yizhu}

\begin{yizhu}[原文疑误]
本片倒数第四个自然段中「取值于 $T$ 的 $(0,1)$-形式 $\alpha_a$」里的 $T$，按前文应读作 $T_C$（$C$ 的全纯切丛）；注意到 $\alpha_a$ 是 $H^1(C,T_C)$ 中元素的代表，原文此处显系笔误，译文照录。
\end{yizhu}

\begin{yizhu}[易错点]
$w=\Tr\,\varphi^2$ 不是函数而是二次微分，它取值于 $K_C^2$，所以 Hitchin 底空间是 ${\EUBB}=H^0(C,K_C^2)$ 而不是某个函数空间。对 $SU(2)$ 有 $\dim {\EUBB}=3g-3=\dim\M=\frac12\dim\MH$；这个「维数正好是一半」的事实正是后面纤维维数公式 (4.25) 的来源，也是完全可积性所需的函数个数。
\end{yizhu}

\begin{yizhu}[含糊点拨]
完全可积性的证明是「先在无穷维空间算、再下降到模空间」的典型：$H_a$ 定义在超凯勒商之前的 $\EUW$ 上，在那里 $\Omega_I$ 使 $\varphi_z$ 与 $A_{\bar z}$ 共轭，而 $H_a$ 只含 $\varphi_z$，故 Poisson 交换是显然的；再由辛约化把 Poisson 交换性带到 $\MH$。按：这也是原文不必写出具体括号计算的原因——真正的内容是「辛约化保持 Poisson 交换性」这一普遍性质。
\end{yizhu}
